\section{LLM-Assisted Assembly Transformation}\label{sec:llm}
\subsection{Task, Representation, and Model Role}
Zone delegates one of EQ, NOP insertion, BCF, or IBP to an LLM,
requesting a complete modified assembly representation. PII and
wrapping remain native transformations. The model generates an
instruction proposal; deterministic preparation, assembly, data
handling, and validation remain part of the evaluated pipeline.

The engine prepares symbolic labels for relevant code locations and
keeps the data representation outside the generated instruction text.
Original bytes, the prepared intermediate representation, and final
output identity are recorded separately. In particular, similarity to
the intermediate input does not measure every change from the original
file. The experiments request the model identifier
\texttt{deepseek-flash}. Version names assigned by a local alias or
pricing table are not independent confirmation of the provider's
model version.

\subsection{Zero-Shot, Two-Shot, and Call Budgets}
The initial zero-shot request supplies a textual specification without
worked examples; the two-shot request adds two input/output examples.
Example identities and prompt versions are part of the condition.
Repeated run numbers do not imply controlled or independent service
seeds. After rejection, the general pipeline can supply the previous
answer and validation feedback, so any later attempt is identified
as feedback-assisted.

The benign benchmark deliberately disables these internal retries:
each invocation allows one generation, and each input/family/shot
cell has three invocations in total, including failures. No failed
invocation is replaced. The ongoing payload campaign instead uses
one target output per cell and a driver supporting up to three
generation attempts. These are different experimental designs and
are not merged into one success percentage.

\subsection{Validation and Independent Functional Assessment}
Structural checks cover assembly, decoding coverage, tracked control
flow, data consistency, and family-specific requirements. A proposal
can be rejected because it is malformed, loses required structure,
or fails to apply the requested transformation. Structural acceptance
alone does not establish behavioral equivalence.

For the new benign benchmark, functional verification follows
generation in a separate harness and compares the complete output
buffer against its predeclared expected contents. It is not a
generation-time selection oracle. This differs from the earlier
MessageBox pilot, where the dialog verifier was included in output
acceptance. That earlier pilot is retained as a separate cohort.

Records preserve accepted and rejected invocations, transport failures,
hashes, timing, usage, and locally estimated cost. Reusing a completed
output is distinguished from a new model call. A correct buffer is
empirical evidence for the fixture's specified observable, not a
proof of all-state equivalence or correctness on unrelated payloads.
